% Cote 119: offprints, reports, courses, CV and letters (typed and printed part): archive cote 119, site page 71 (PDF page 72).
% offprint. Draft compared: none.
% AI-transcribed by Claude Opus 5.5 (claude-opus-5-5, adaptive thinking, effort high), Message Batches API, batch r119_full (msgbatch_01DwZNm5Y7YsKBWyAN4sFs8w), 2026-11-13; not yet checked by eye.
\scanpage{119/71}
\mspage{( 3 )}

De même que dans 1.2, on a un morphisme canonique de faisceaux
en monoïdes
\[ (2.1.1) \qquad G \to End^{\otimes}(\omega_0^G). \]

Proposition 2.2. --- Soient $k$ un corps, G un $k$-monoïde affine, T un
$k$-schéma. On a alors :

($a$) Le morphisme (2.1.1) est un isomorphisme.

($b$) Si $\omega$, $\omega'$ sont des foncteurs fibre sur $\mathrm{Rep}_0(G)$ à valeurs dans T, et
si G est un groupe, $\omega$ et $\omega'$ sont localement isomorphes pour la topologie fpqc.

La collection des catégories $\mathrm{Fib}_0(G, T)$ définit de façon naturelle une
catégorie fibrée [(\unsure{${}^6$}), VI, 6.1] au-dessus de $\mathrm{Sch}_{/k}$, qu'on notera $\mathrm{FIB}_0(G)$.
Si TORS(G) dénote la catégorie fibrée au-dessus de $\mathrm{Sch}_{/k}$ des torseurs
sous $G_T$ pour T variable [(\unsure{${}^5$}), loc. cit.], la proposition 4.2 peut encore
s'énoncer comme suit :

Corollaire 2.3. --- Le foncteur $\omega \mapsto \mathrm{Isom}^{\otimes}(\omega_0^G, \omega)$ définit une équi-%
valence de catégories fibrées
\[ (2.3.1) \qquad \mathrm{FIB}_0(G) \xrightarrow{\sim} \mathrm{TORS}(G) \]
au-dessus de $\mathrm{Sch}_{/k}$.

La partie ($a$) de 2.2 avait été signalée par Cartier [(${}^3$), th. 1] dans le
cas d'un groupe algébrique linéaire.

3. Un théorème d'existence.

3.1. On peut donner une caractérisation axiomatique des catégories
Rep(G) pour G un S-monoïde affine, et S un schéma affine. Si S est le
spectre d'un corps $k$, celle-ci prend la forme suivante :

Théorème 3.2. --- Soit $C$ une $\otimes$-catégorie ACU $k$-linéaire. Alors $C$ est
$\otimes$-équivalente à une $\otimes$-catégorie $\mathrm{Rep}_0(G)$ avec G un $k$-monoïde affine si
et seulement s'il existe un $\otimes$-foncteur ACU $k$-linéaire, exact et fidèle
$\omega : C \to \mathrm{Modf}(k)$. Le $k$-monoïde G est un groupe si et seulement si $C$ possède
des objets Hom (${}^7$).

4. Quelques traductions. --- La correspondance établie plus haut
entre groupes et $\otimes$-catégories permet de traduire des propriétés de théorie
des groupes dans un langage catégorique. Avant d'en donner quelques
exemples, posons une définition. Si V est un objet d'une $\otimes$-catégorie ACU
abélienne $C$, on dit que V est un $\otimes$-générateur de $C$, si tout objet de $C$
est quotient d'un sous-objet d'un objet de la forme P(V), $P \in \mathbf{N}[t]$, où
dans l'expression P(V), la somme (resp. le produit) est remplacé par $\oplus$
(resp. $\otimes$).

Soient G, G$'$ des $k$-groupes ($k$ un corps), $u : G \to G'$ un morphisme
de $k$-groupes, V un G-$k$-module de rang fini sur $k$.

\mspage{\hand{71}}
